%% LyX 2.4.4 created this file.  For more info, see https://www.lyx.org/.
%% Do not edit unless you really know what you are doing.
\documentclass[english]{article}
\usepackage[T1]{fontenc}
\usepackage[latin9]{inputenc}
\usepackage{enumitem}
\usepackage{amsmath}
\usepackage{amsthm}
\usepackage{amssymb}
\usepackage{geometry}
\geometry{verbose,tmargin=1in,bmargin=1in,lmargin=1in,rmargin=1in}
\usepackage{wasysym}

\makeatletter

%%%%%%%%%%%%%%%%%%%%%%%%%%%%%% LyX specific LaTeX commands.
%% Because html converters don't know tabularnewline
\providecommand{\tabularnewline}{\\}

%%%%%%%%%%%%%%%%%%%%%%%%%%%%%% Textclass specific LaTeX commands.
\numberwithin{equation}{section}
\numberwithin{figure}{section}
\newlength{\lyxlabelwidth}      % auxiliary length 

%%%%%%%%%%%%%%%%%%%%%%%%%%%%%% User specified LaTeX commands.
\usepackage{fancyhdr}
\usepackage{lastpage}
\pagestyle{fancy}
\usepackage{enumitem}
\usepackage{ifthen}
%sets math fonts to be more readable
\everymath=\expandafter{\the\everymath\displaystyle}
%\DeclareMathSizes{10}{10}{10}{10}
\usepackage{multicol}

\usepackage{advdate}    % Advancing/saving dates
\usepackage[dayofweek]{datetime}   % Dates formatting
\usepackage{datenumber} % Counters for dates
% Added by lyx2lyx
\setlength{\parskip}{\smallskipamount}
\setlength{\parindent}{0pt}

\makeatother

\usepackage{babel}
\begin{document}
\renewcommand{\author}{Dr.\,James Ashe}

\newcommand{\classNum}{335}

\newcommand{\classSec}{A}

\newcommand{\nameType}{Name:}

\newcommand{\examType}{Homework 3}\usdate

\newcommand{\Date}{Due date: \formatdate{18}{3}{2013}} %%\AdvanceDate[12] \today

%\newdate{my_date}{\day}{\month}{\year}%   
%\AdvanceDate[-5] 
%\displaydate{my_date}

%%First page's header%%%%%%%%%%%%%%%%%%%%%%
\fancypagestyle{firststyle} %%header settings for first pagr
{
	\fancyhead[L]{MTH \classNum{}(\classSec{}) \author{}\\
	\textbf{\examType{}} \Date{}}
	\fancyhead[C]{\textbf{\nameType}}
	\setlength{\headsep}{14pt}
}
%%%%%%%%%%%%%%%%%%%%%%%%%%%%%%%%%%%%%%%%%%

%%Next page's header%%%%%%%%%%%%%%%%%%%%%%%
%\setlist{nolistsep} %eliminates space between list items
\lhead{MTH \classNum{}(\classSec{})} 
\chead{\textbf{\nameType}} 
\cfoot{\thepage{}~of~\pageref{LastPage}} 
\setlength{\headsep}{5pt} 
%%%%%%%%%%%%%%%%%%%%%%%%%%%%%%%%%%%%%%%%%%%

%%Various settings; see the preamble for other settings%%%%%
%%\setenumerate[0]{leftmargin=12pt,itemindent=0pt,labelwidth=0pt}
\renewcommand{\labelenumi}{\textbf{\arabic{enumi}.}}
\renewcommand{\labelenumii}{\textbf{\alph{enumii}.}}
\thispagestyle{firststyle}
%%%%%%%%%%%%%%%%%%%%%%%%%%%%%%%%%%%%%%%%%%

\textbf{Homework Problems. }\emph{Solutions should be written/typed
in numerical order and clearly labeled. Unsolved problems should be
included but left blank.}

Use the following to answer the next two problems. Let $E=\{1,2,3\}$.
Consider the following relations in $E$.
\begin{enumerate}[resume]
\item State whether each relation below, $\mathcal{R}_{1},\mathcal{R}_{2},\mathcal{R}_{3},\mathcal{R}_{4},\mathcal{R}_{5}$,
is symmetric. Justify your answers.
\end{enumerate}
\begin{tabular}{ll}
$\mathcal{R}_{1}=\left\{ \left(1,1\right),\left(2,1\right),\left(2,2\right),\left(3,2\right),\left(2,3)\right)\right\} $ & $\mathcal{R}_{4}=\left\{ \left(1,1\right),\left(3,2\right),\left(2,3\right)\right\} $\tabularnewline
$\mathcal{R}_{2}=\left\{ \left(1,1\right)\right\} $ & $\mathcal{R}_{5}=E\times E$\tabularnewline
$\mathcal{R}_{3}=\left\{ \left(1,2\right)\right\} $ & \tabularnewline
\end{tabular}\vfill
\begin{enumerate}[resume]
\item State whether each relation below, $\mathcal{R}_{6},\mathcal{R}_{7},\mathcal{R}_{8},\mathcal{R}_{9},\mathcal{R}_{10}$,
is transitive. Justify your answers.
\end{enumerate}
\begin{tabular}{ll}
$\mathcal{R}_{6}=\left\{ \left(1,2\right),\left(2,2\right)\right\} $ & $\mathcal{R}_{9}=\left\{ \left(1,1\right)\right\} $\tabularnewline
$\mathcal{R}_{7}=\left\{ \left(1,2\right),\left(2,3\right),\left(1,3\right),\left(2,1\right),\left(1,1\right)\right\} $ & $\mathcal{R}_{10}=E\times E$\tabularnewline
$\mathcal{R}_{8}=\left\{ \left(1,2\right)\right\} $ & \tabularnewline
\end{tabular}\vfill
\begin{enumerate}[resume]
\item Define the following relation on the real numbers, $\mathcal{R}:x\mathcal{R}y\iff xy\geq0$.
Prove that $\mathcal{R}$ is reflexive and symmetric, but that $\mathcal{R}$
is not an equivalence relation.\vfill
\item Define the following equivalence relation on the set of subsets of
$X$ as follows: $A\mathcal{R}B\iff A\cap B\neq\emptyset$ for $A,B\subseteq X$
Prove or disprove that $\mathcal{R}$ is an equivalence relation.
\vfill
\item Let $\mathcal{R}$ and $\mathcal{R}'$ be symmetric relations in a
set $A$. Prove that $\mathcal{R}\cap\mathcal{R}'$ is a symmetric
relation in A. (note that $\mathcal{R}$ and $\mathcal{R}'$ are subsets
of $A\times A$). \vfill
\item Determine the following:

\begin{enumerate}[resume]
\item $\gcd(2^{4}3^{2}5\cdot7^{2},2\cdot3^{3}7\cdot11)$
\item The lowest common multiple: $\mbox{lcm}(2^{4}3^{2}5\cdot7^{2},2\cdot3^{3}7\cdot11)$.\vfill
\end{enumerate}
\item Determine $51\mod 13$, $342\mod 85$, and $(82\cdot73)\mod 7$.\vfill
\item Find integers $s$ and $t$ so that $1=7\cdot s+11\cdot t$. Show
that $s$ and $t$ are not unique. \vfill
\item Let $n$ be a positive integer greater than $1$ and let $a$ and
$b$ be positive integers. Prove that $a\mod n=b\mod n$ if and only
if $a=b\mod n$. \vfill
\item Let $d=\gcd(a,b)$. If $a=da'$ and $b=db'$, show that $\gcd(a',b')=1$.\vfill\newpage
\end{enumerate}
\textbf{Presentation problems.}
\begin{enumerate}
\item Let the function $f:A\rightarrow B$ be 1-1 and onto, i.e., $f^{-1}$
exists. Prove that $\left(f^{-1}\circ f\right):A\rightarrow A$ is
the identity function on $A$ and that $\left(f\circ f^{-1}\right):B\rightarrow B$
is the identity function on $B$.
\item Let $f:A\rightarrow B$ be 1-1 and onto. Prove that $g=f^{-1}$, if
$\left(g\circ f\right):A\rightarrow A$ is the identity function on
$A$ and $\left(f\circ g\right):B\rightarrow B$ is the identity function
on $B$. 
\item The \emph{product set }of sets $A$ and $B$ is the set of all ordered
pairs $(a,b)$ where $a\in A$ and $b\in B$; it is denoted by $A\times B$.
Prove that $A\times(B\cap C)=(A\times B)\cap(A\times C)$. 
\item Prove: Let $R$ and $R'$ be symmetric relations in a set $A$; then
$R\cap R'$ is a symmetric relation in $A$.
\item Prove that ${\displaystyle \sum_{k=1}^{n}2k=n^{2}+n}$
\item The Fibonacci numbers are :$1,1,2,3,5,8,13,21,34,\dots$. In general,
the Fibonacci numbers are defined by $f_{1}=1$, $f_{2}=1$, and $f_{n+2}=f_{n+1}+f_{n}$
for $n=1,2,3,\dots$. Prove that the $n^{\mbox{th}}$ Fibonacci number
$f_{n}<2^{n}$.
\item Prove that there infinitely many prime numbers. 
\item [$\smiley$\textbf{.}]Two sets are said to be \emph{equal }if there
exists a bijective function between them; we write $\left|A\right|=\left|B\right|$.
$B$ is said to have \emph{strictly greater cardinality} than $B$
if there is a one-to-one function from $A$ into $B$, but there does
not exist a bijective function from $A$ into $B$; we write$\left|A\right|<\left|B\right|$.
Prove that $\left|\mathbb{N}\right|<\left|\mathbb{R}\right|$.
\end{enumerate}

\end{document}
